\section{Introduction}
\label{sec:intro}

Segmentation maps of overhead imagery routinely mix feature types with fundamentally different
geometry: \emph{linear} features such as roads, one-dimensional structures extruded to a small
width; \emph{polygonal} features such as building footprints, genuinely two-dimensional regions;
and optionally \emph{point} features --- trees, lamp posts, manhole covers --- small discrete
objects spanning only a handful of pixels. A single evaluation tile covering a city block can
legitimately contain all three. Yet the standard tool for scoring segmentation quality, pixel IoU,
treats every class identically as a region-overlap problem, and that assumption breaks down
specifically for the linear case: because a road is only a few pixels wide, IoU denominators are
dominated by exactly the pixels most sensitive to sub-pixel-scale registration noise, centerline
jitter, and predicted-width error, none of which reflect whether the road network was actually
extracted correctly. The same problem recurs one level up: the \emph{training} losses
conventionally used for multiclass segmentation (cross-entropy, soft Dice) are exactly as
geometry-agnostic as IoU is as an evaluation tool --- they have no notion of tolerance, centerline
topology, or point-instance identity either, so a model trained with them has no direct pressure to
get those properties right, only to match pixels.

Table~\ref{tab:failure-modes} previews five concrete scenarios on a controlled $128\times128$
synthetic scene (a 3px-wide vertical road stripe plus a $40\times40$ building rectangle; full
setup in Section~\ref{sec:motivation}), with values computed directly against the live metrics
library for this paper. It shows where IoU's verdict diverges from actual prediction quality, and
--- importantly --- a case where \textbf{DTAF1 itself is fooled} even though it was designed to
fix IoU's problem: at 33\% random road-pixel deletion, DTAF1 still reports a perfect
$\mathbf{1.000}$, motivating DTAF1-Topo in Section~\ref{sec:topological}.

\begin{table*}
  \caption{IoU and DTAF1 failure modes for linear and polygonal features.}
  \label{tab:failure-modes}
  \begin{tabular}{@{}lccccc l@{}}
    \toprule
    Scenario & Road IoU & clDice & DTAF1 & CBHM & CBHM-soft & Correct read \\
    \midrule
    Perfect prediction & 1.000 & 1.000 & 1.000 & 1.000 & 1.000 & baseline \\
    Road shifted 5px & 0.000 & 0.000 & \textbf{1.000} & 0.000 & 0.806 &
      centerline correct; IoU/CBHM wrongly fail it \\
    Road $3.3\times$ too thick (10px vs.\ 3px GT) & 0.281 & 1.000 & 1.000 & 1.000 & 1.000 &
      centerline correct, only width wrong \\
    Road 33\% pixels deleted & 0.667 & 0.759 & \textbf{1.000} & 0.863 & 0.953 &
      network degraded; \textbf{DTAF1 misses this} \\
    Building eroded 5px & $0.562^\dagger$ & --- & 0.917 & 0.939 & 0.907 &
      shape genuinely worse; both degrade correctly \\
    Sparse road (15px offset, ${\sim}0.4\%$ of image) & 0.000 & --- & 0.500 & \textbf{0.000} & 1.000 &
      building perfect, road wrong; CBHM collapse \\
    Null prediction & 0.000 & 0.000 & 0.000 & 0.000 & 0.000 & baseline \\
    \bottomrule
  \end{tabular}
  \vspace{2pt}
  {\footnotesize $^\dagger$Building IoU, not road IoU, for this row.}
\end{table*}

We adopt a simple, uniform data representation throughout: every input is an $H\times W$
\texttt{uint8} NumPy array of integer class labels, with $0=$ background, $1=$ road, $2=$
building, and an optional $3=$ point feature, used only where the underlying dataset actually has
a point class. Every distance-based tolerance in the library is expressed in pixels but is meant
to be tied to a tile's ground sample distance (GSD): $d = \text{physical\_metres} /
\text{GSD\_metres\_per\_pixel}$, so the same physical tolerance (e.g., ``3 metres of road
centerline slop'') applies consistently across imagery captured at different resolutions.
Section~\ref{sec:method} reuses this exact class convention on batched torch tensors --- channel
index equals class id --- so a \texttt{class\_config} dict is literally shareable between an eval
call (\texttt{dtaf1()}) and a training loss (\texttt{ToleranceBandLoss}).

This paper makes seven contributions:

\begin{enumerate}
  \item \textbf{A differentiable multi-task training loss} (Section~\ref{sec:method}) whose three
    geometric terms are direct, well-precedented relaxations of three of this paper's own eval
    metrics --- soft tolerance-band (DTAF1), soft-clDice (Shit et~al.\ \cite{shit2021cldice}), and
    Gaussian-heatmap point regression (CenterNet-style) --- plus a masked multi-source training
    scheme letting heterogeneous datasets (SpaceNet road+building, Potsdam building+point) jointly
    supervise one 4-class model despite neither alone annotating all three feature types.
    Validated via controlled synthetic ablations (Section~\ref{sec:results}) and a small real
    joint-training sanity check across 153 real tiles (Section~\ref{sec:results}) --- not a
    benchmark-scale study.
  \item \textbf{DTAF1} (Section~\ref{sec:evalprotocol}) --- a single, class-agnostic
    tolerance-radius F1 formula spanning linear and polygonal classes with one code path; scoring
    a newly added feature type is a configuration change, not a code change.
  \item \textbf{CBHM} (Section~\ref{sec:evalprotocol}) --- a harmonic-mean composite of clDice and
    boundary F1, deliberately designed as DTAF1's ``harsh'' foil, so a single feature type failing
    outright is not masked by the other type succeeding.
  \item \textbf{\texttt{dtaf1\_weighted} / \texttt{cbhm\_soft}} (Section~\ref{sec:composite}) ---
    GT-pixel-count-weighted companions to DTAF1 and CBHM that soften single-class collapse,
    motivated by a real failure case found on SpaceNet imagery (\texttt{Khartoum\_img371}).
  \item A raster-only approximation of \textbf{APLS} (Section~\ref{sec:evalprotocol}) --- no
    vector road graph survives anywhere in this pipeline, see Section~\ref{sec:implementation} ---
    and \textbf{DTAF1-Topo} (Section~\ref{sec:topological}), which blends this connectivity term
    into DTAF1 and closes the road-breakage blind spot demonstrated in
    Table~\ref{tab:failure-modes}.
  \item \textbf{\texttt{point\_f1}} (Section~\ref{sec:evalprotocol}), a point-feature metric using
    Hungarian instance matching rather than greedy nearest-neighbor assignment, avoiding a
    specific adversarial failure mode.
  \item Empirical validation (Sections~\ref{sec:expsetup}--\ref{sec:results}) spanning nine
    synthetic eval-metric sensitivity sweeps, seven synthetic loss-ablation sweeps, a curated
    real-data sample from SpaceNet (four cities) and ISPRS Potsdam, and a first trained-model
    pipeline scored end-to-end by the library.
\end{enumerate}

The rest of the paper is organized as follows. Section~\ref{sec:related} places this work relative
to prior evaluation metrics and differentiable segmentation losses. Section~\ref{sec:evalprotocol}
defines each eval metric primitive, following a one-subsection-per-primitive structure.
Section~\ref{sec:motivation} empirically motivates why these primitives are necessary.
Sections~\ref{sec:composite} and~\ref{sec:topological} present two analysis-and-design studies ---
composite robustness and topological robustness. Section~\ref{sec:method} introduces the
differentiable multi-task training loss, presenting each term as a direct analog of one primitive
from Section~\ref{sec:evalprotocol}. Section~\ref{sec:implementation} describes the
implementation. Sections~\ref{sec:expsetup}--\ref{sec:results} give the full experimental setup
and results across synthetic eval sweeps, synthetic loss ablations, real-data, and trained-model
settings. Section~\ref{sec:discussion} discusses limitations honestly. Section~\ref{sec:conclusion}
concludes.
